\Introduction

В современных образовательных учреждениях управление расписанием учебных занятий является одной из ключевых организационных задач. Диспетчеру по расписанию необходимо оперировать большим объёмом данных: учебные группы, преподаватели, аудитории, дисциплины, семестры---и обеспечивать их непротиворечивую связку. Традиционные подходы (Excel-таблицы, бумажные журналы) приводят к ошибкам, дублированию информации и низкой оперативности внесения изменений.

Автоматизированное рабочее место (АРМ) диспетчера по расписанию призвано решить задачи централизованного учёта всех элементов расписания, импорта данных из Excel-файлов, визуализации недельной сетки занятий и экспорта расписания в формат iCalendar для интеграции с календарями Google и Apple.

Целью настоящей работы является разработка АРМ диспетчера по расписанию---клиент-серверного приложения с REST API, десктопным клиентом на C++/Qt6 и веб-интерфейсом. Для достижения поставленной цели необходимо решить следующие задачи:

\begin{itemize}
\item проектирование схемы базы данных и REST API на основе OpenAPI 3.0;
\item реализация серверной части на Go с автогенерацией кода по OpenAPI-спецификации;
\item разработка десктопного клиента на C++/Qt6 с автогенерацией API-клиента;
\item реализация импорта расписания из ZIP-архивов с XLSX-файлами;
\item реализация экспорта расписания в формат iCalendar;
\item контейнеризация приложения средствами Docker.
\end{itemize}
